% !TeX program = xelatex
% ============================================================================
%  第 4 课　群的四条规则　—— 学生版讲义
%  与 02-课件/第04课-群的四条规则/第04课-群的四条规则.tex 同步
% ============================================================================

\xhlesson{4}{群的四条规则}{}

\begin{mubiao}
  \begin{itemize}
    \item 能说出群的四条规则。
    \item 会在模 \(12\) 加法、有理数加法、新定义运算里各找到一个例子。
    \item 知道单位元不一定是 \(0\)，交换律不是群的必要条件。
  \end{itemize}
\end{mubiao}

\section*{一、运算}

\begin{definition}[运算]
  设 \(S\) 是一个集合。如果有一种办法，把 \(S\) 里任意两个元素 \(a,b\) 合起来，
  得到 \(S\) 里\textbf{唯一确定}的一个元素，记作 \(a*b\)，
  就说 \(*\) 是 \(S\) 上的一个\textbf{运算}。
\end{definition}

\begin{sikao}
  「取两个数里较大的那一个」是运算吗？试着举一个小例子说说看。
  \liukong{2}
\end{sikao}

\section*{二、群：四条规则}

\begin{definition}[群]
  设 \(*\) 是集合 \(G\) 上的一个运算。如果下面四条都成立，
  就说 \((G,*)\) 是一个\textbf{群}：
  \begin{enumerate}
    \item \textbf{封闭}：任取 \(a,b\in G\)，\(a*b\) 仍然属于 \(G\)；
    \item \textbf{结合}：\((a*b)*c=a*(b*c)\)；
    \item \textbf{单位元}：存在 \(e\in G\)，对每个 \(a\) 都有 \(a*e=e*a=a\)；
    \item \textbf{逆元}：每个 \(a\in G\) 都有 \(b\in G\)，使 \(a*b=b*a=e\)。
  \end{enumerate}
\end{definition}

\begin{sikao}
  回头看看上节课那张表，数一数：上面写了几条规则？\underline{\hspace{2em}} 条。\\
  表的标题里，有没有「交换律」？\underline{\hspace{2em}}
\end{sikao}

\section*{三、四个例子，逐条打勾}

\begin{example}[先看三个老熟人]
  \begin{center}
  \begin{tabular}{@{}lcccc@{}}
    \toprule
     & 封闭 & 结合 & 单位元 & 逆元 \\
    \midrule
    模 \(12\) 加法 & & & & \\
    有理数加法 & & & & \\
    有理数乘法 & & & & \\
    \bottomrule
  \end{tabular}
  \end{center}
  在每格里打 \(\checkmark\) 或 \(\times\)，并把单位元、逆元写出来。
\end{example}

\begin{sikao}
  有理数乘法里，\(0\) 的逆元是什么？
  \liukong{2}
  所以有理数乘法是不是群？\underline{\hspace{4em}}
\end{sikao}

\begin{example}[自己定制一个运算]
  规定 \(a\oplus b=a+b+1\)。算一算：
  \[
    3\oplus5=\underline{\hspace{3em}}, \qquad (-2)\oplus0=\underline{\hspace{3em}} .
  \]
\end{example}

\begin{dongshou}[一条一条验]
  \begin{itemize}
    \item 封闭：\(a\oplus b\) 还是有理数吗？
    \item 结合：\((a\oplus b)\oplus c=\underline{\hspace{5em}}\)，
      \(a\oplus(b\oplus c)=\underline{\hspace{5em}}\)，两边一样吗？
    \item 单位元：解 \(a+e+1=a\)，得 \(e=\underline{\hspace{3em}}\)。是 \(0\) 吗？
    \item 逆元：解 \(a+b+1=-1\)，得 \(b=\underline{\hspace{3em}}\)。
  \end{itemize}
  \liukong{3}
\end{dongshou}

\begin{dongshou}[用手指回答]
  在 \(a\oplus b=a+b+1\) 里，报出下面每个数的逆元：
  \begin{center}
  \begin{tabular}{@{}lccc@{}}
    \toprule
    \(a\) & \(0\) & \(1\) & \(-3\) \\
    \midrule
    \(a\) 的逆元 & \hspace{1.6em} & \hspace{1.6em} & \hspace{1.6em} \\
    \bottomrule
  \end{tabular}
  \end{center}
\end{dongshou}

\section*{四、结合律不是废话}

\begin{dongshou}
  再换一个运算：\(a\oplus b=\dfrac{a+b}{2}\)。算一算：
  \[
    (1\oplus3)\oplus5=\underline{\hspace{4em}}, \qquad
    1\oplus(3\oplus5)=\underline{\hspace{4em}} .
  \]
  两个结果一样吗？\underline{\hspace{4em}}\quad 所以它满足结合律吗？\underline{\hspace{4em}}
\end{dongshou}

\section*{五、交换律不是必要的}

\begin{example}
  动作甲：把数\textbf{乘 \(2\)}。\qquad 动作乙：把数\textbf{加 \(1\)}。

  取 \(x=1\)，两条路：

  \begin{itemize}
    \item 先甲后乙：\(1\to\underline{\hspace{2em}}\to\underline{\hspace{2em}}\)；
    \item 先乙后甲：\(1\to\underline{\hspace{2em}}\to\underline{\hspace{2em}}\)。
  \end{itemize}

  结果一样吗？\underline{\hspace{4em}}
\end{example}

\begin{dongshou}
  三个动作连着做：甲、乙、甲。

  先合前两个：（甲，乙）再甲；\qquad 先合后两个：甲 再（乙，甲）。

  我猜结果\underline{\hspace{2em}}（一样／不一样）。
  \liukong{2}
\end{dongshou}

\begin{xiaojie}
  用你的话说说：为什么说「交换律不是群的必要条件」？
  \liukong{4}
\end{xiaojie}

\noindent\textbf{下节课}：三角形的对称——正三角形有 \(6\) 个对称动作，我们把它们写成两行式。